\chapter{اللقاء الرابع: نزول القرآن على سبعة أحرف}

\section{القول في اللغة التي نزل بها القرآن}
\isnad{قال أبو جعفر رحمه الله:}
«قد دللنا على صحة القول بما فيه الكفاية لمن وفقه الله لفهمه، على أن الله جل ثناؤه أنزل جميع القرآن بلسان العرب دون غيرها من ألسن سائر أجناس الأمم، وعلى فساد قول من زعم أن منه ما ليس بلسان العرب ولغتها. فنقول الآن إذ كان ذلك صحيحاً: أبألسن جميعها أم بألسن بعضها؟ إذ كانت العرب وإن جمع جميعها اسم أنهم عرب، وهم مختلفو الألسن بالبيان، ومتباينو المنطق والكلام.

وإن كان ذلك كذلك، وكان الله جل ذكره قد أخبر عباده أنه قد جعل القرآن عربياً، وأنه أنزله بلسان عربي مبين، ثم كان ظاهره محتملاً خصوصاً وعموماً، لم يكن لنا السبيل إلى العلم بما عنى الله تعالى ذكره من خصوصه وعمومه إلا ببيان من جُعل إليه بيان القرآن، وهو رسول الله صلى الله عليه وسلم.

وإذ كان ذلك كذلك، وكانت الأخبار قد تظاهرت عنه صلى الله عليه وسلم بما حدثنا به...».

\begin{fawaid}[منهج الطبري في عرض المسائل]
نلاحظ هنا منهج الإمام \rawi{الطبري} في عرض المسائل الكبرى؛ فهو يبدأ أولاً بجمع كل ما ورد في الباب من روايات وآثار، فيسوقها سرداً دون تعليق على صحتها أو ضعفها في الغالب، ثم بعد أن يجمع المادة العلمية كاملة، يبدأ في تحليلها ونقدها وبيان رأيه فيها. وهذا منهج دقيق يدل على استيعابه وحرصه على بناء قوله على استقراء تام.
\end{fawaid}

\section{الأحاديث الواردة في نزول القرآن على سبعة أحرف}
ساق \rawi{الطبري} رحمه الله عدداً كبيراً من الأحاديث والآثار في هذا الباب، نذكر منها ما يلي على سبيل المثال لا الحصر:

\begin{itemize}
    \item \isnad{عن} \rawi{أبي هريرة} \isnad{أن رسول الله صلى الله عليه وسلم قال:} \begin{hadith}أُنزل القرآن على سبعة أحرف، فالمراء في القرآن كفر (ثلاث مرات)، فما عرفتم منه فاعملوا به، وما جهلتم منه فردوه إلى عالمه\end{hadith}.
    \item \isnad{عن} \rawi{عبدالله بن مسعود} \isnad{قال: قال رسول الله صلى الله عليه وسلم:} \begin{hadith}أُنزل القرآن على سبعة أحرف، لكل حرف منها ظهر وبطن، ولكل حرف حد، ولكل حد مطلع\end{hadith}.
    \item \isnad{عن} \rawi{عمر بن الخطاب} \isnad{أنه قال:} \begin{hadith}سمعت \rawi{هشام بن حكيم} يقرأ سورة الفرقان في حياة رسول الله صلى الله عليه وسلم، فاستمعت لقراءته، فإذا هو يقرؤها على حروف كثيرة لم يقرئنيها رسول الله صلى الله عليه وسلم كذلك، فكدت أساوره في الصلاة، فتصبرت حتى سلم، فلببته بردائه فقلت: من أقرأك هذه السورة؟ قال: أقرأنيها رسول الله صلى الله عليه وسلم. فقلت: كذبت، فوالله إن رسول الله صلى الله عليه وسلم هو أقرأني هذه السورة التي سمعتك. فانطلقت به أقوده إلى رسول الله صلى الله عليه وسلم، فقلت: يا رسول الله، إني سمعت هذا يقرأ سورة الفرقان على حروف لم تقرئنيها، وأنت أقرأتني سورة الفرقان. فقال رسول الله صلى الله عليه وسلم: أرسله يا عمر، اقرأ يا هشام. فقرأ عليه القراءة التي سمعته يقرؤها، فقال رسول الله صلى الله عليه وسلم: هكذا أنزلت. ثم قال: اقرأ يا عمر. فقرأت القراءة التي أقرأني، فقال: هكذا أنزلت. إن هذا القرآن أنزل على سبعة أحرف، فاقرؤوا ما تيسر منه\end{hadith}.
    \item \isnad{عن} \rawi{أبي بن كعب} \isnad{أن رسول الله صلى الله عليه وسلم قال:} \begin{hadith}أقرأني جبريل على حرف، فراجعته، فلم أزل أستزيده فيزيدني حتى انتهى إلى سبعة أحرف\end{hadith}. \isnad{قال} \rawi{ابن شهاب}: «بلغني أن تلك السبعة الأحرف إنما هي في الأمر الذي يكون واحداً، لا يختلف في حلال ولا حرام».
    \item \isnad{عن} \rawi{أبي بكرة} \isnad{عن أبيه قال: قال رسول الله صلى الله عليه وسلم:} \begin{hadith}قال جبريل: اقرأ القرآن على حرف. قال ميكائيل: استزده. حتى بلغ سبعة أحرف، قال: كلها شافٍ كافٍ، ما لم تختم آية عذاب بآية رحمة، أو آية رحمة بآية عذاب، كقولك: هلم وتعال\end{hadith}.
\end{itemize}

\separator

\section{تحليل الطبري لمعنى الأحرف السبعة}
بعد أن ساق \rawi{الطبري} الروايات، بدأ في تحليلها ليصل إلى معنى «الأحرف السبعة».

\isnad{قال أبو جعفر:} «صح وثبت أن الذي نزل به القرآن من ألسن العرب البعض منها دون الجميع، إذ كان معلوماً أن ألسنتها ولغاتها أكثر من سبعة».

ثم طرح سؤالاً: ما البرهان على أن معنى الحديث هو سبع لغات، وليس سبعة أوجه من المعاني (كالأمر والزجر والوعد والوعيد) كما قال آخرون؟

\subsection{حجة الطبري في مذهبه}
يرى \rawi{الطبري} أن سياق الأحاديث نفسها هو الحجة الدامغة على أن الاختلاف كان في \textbf{اللفظ والتلاوة} لا في \textbf{المعنى والحكم}.

\isnad{قال رحمه الله:} «إنهم تماروا في القرآن فخالف بعضهم بعضاً في نفس التلاوة دون ما في ذلك من المعاني، وإنهم احتكموا فيه إلى النبي صلى الله عليه وسلم، فاستقرأ كل رجل منهم، ثم صوب جميعهم في قراءتهم على اختلافها... فلو كان اختلافهم فيما دلت عليه تلاوتهم من التحليل والتحريم والوعد والوعيد، لكان مستحيلاً أن يصوبهم جميعاً رسول الله صلى الله عليه وسلم... لأن ذلك لو جاز، لوجب أن يكون الله قد أمر بشيء ونهى عنه في نفس الوقت، وذلك إثبات لما قد نفى الله عن كتابه فقال: \begin{ayah}أَفَلَا يَتَدَبَّرُونَ الْقُرْآنَ ۚ وَلَوْ كَانَ مِنْ عِندِ غَيْرِ اللَّهِ لَوَجَدُوا فِيهِ اخْتِلَافًا كَثِيرًا\end{ayah}».

ثم استدل بالرواية التي فيها: \begin{hadith}كقولك هلم وتعال\end{hadith}.

\isnad{قال:} «فقد أوضح نص هذا الخبر أن اختلاف الأحرف السبعة إنما هو اختلاف ألفاظ، كقولك: هلم، وتعال، باتفاق المعاني، لا باختلاف معانٍ موجبة اختلاف أحكام».

\begin{fawaid}[خلاصة مذهب الطبري في الأحرف السبعة]
يرى \rawi{الطبري} أن الأحرف السبعة هي سبع لغات (ألسن) من لغات العرب، نزل بها القرآن على وجه التخيير والتيسير، وهي متفقة في المعنى وإن اختلفت في اللفظ، كالمترادفات (هلم، تعال، أقبل). وهو بذلك ينفي أن يكون المراد بها سبعة أبواب من المعاني (كالأمر والنهي)، ويرى أن هذا الفهم الأخير هو تفسير لحديث آخر وهو: «أنزل القرآن من سبعة أبواب من الجنة».
\end{fawaid}

\section{إشكالات وردود حول مذهب الطبري}
بعد أن قرر \rawi{الطبري} مذهبه، استعرض الاعتراضات المحتملة عليه وأجاب عنها، وهذا من قوة منهجه.

\subsection{الإشكال الأول: أين نجد اليوم حرفاً مقروءاً بسبع لغات؟}
اعترض عليه معترض فقال: «إن كان الأمر كما تقول، فأوجدنا حرفاً واحداً في كتاب الله مقروءاً بسبع لغات لنسلم لك».

\isnad{فأجاب الطبري:} «إنا لم ندَّعِ أن ذلك موجود اليوم، وإنما أخبرنا أن معنى قول النبي صلى الله عليه وسلم... هو ما وصفنا».

\subsection{الإشكال الثاني: ما مصير الأحرف الستة الباقية؟}
اعترض عليه معترض آخر: «فما بال الأحرف الستة غير موجودة اليوم؟ هل نُسخت ورُفعت؟ أم ضيعتها الأمة؟».

\isnad{فأجاب الطبري:} «لم تُنسخ فتُرفع، ولا ضيعتها الأمة، ولكن الأمة أُمرت بحفظ القرآن وخُيّرت في قراءته بأي تلك الأحرف السبعة شاءت... فلما رأت أن الاختلاف والتنازع قد أفضى ببعضهم إلى التكذيب والتكفير، اجتمعوا على حرف واحد نظراً منهم لأنفسهم ولمن بعدهم، وتركوا القراءة بالأحرف الستة الباقية طاعةً لإمامهم العادل (\rawi{عثمان بن عفان}) ورأفةً بالأمة».

واستشهد على ذلك بقصة جمع القرآن في عهد \rawi{أبي بكر} بعد موقعة اليمامة، ثم قصة جمع \rawi{عثمان} للناس على مصحف واحد بعد اختلاف أهل الشام والعراق في القراءة في غزوة أرمينية.

\subsection{الإشكال الثالث: كيف جاز للصحابة ترك قراءة أقرأهم إياها النبي؟}
\isnad{أجاب الطبري:} «إن أمره إياهم بذلك لم يكن أمر إيجاب وفرض، وإنما كان أمر إباحة ورخصة... فلم يكن القوم بتركهم نقل جميع القراءات السبع تاركين ما كان عليهم نقله، بل كان الواجب عليهم من الفعل ما فعلوا، إذ كان الذي فعلوا من ذلك هو النظر للإسلام وأهله».

\section{الفرق بين الأحرف السبعة والقراءات}
نبه \rawi{الطبري} على مسألة دقيقة جداً، وهي أن الخلاف في القراءات المعروف اليوم (في رفع حرف أو نصبه أو جره) ليس هو المقصود بحديث الأحرف السبعة.

\isnad{قال:} «فأما ما كان من اختلاف القراءة في رفع حرف وجره ونصبه، وتسكين حرف وتحريكه... فمن معنى قول النبي صلى الله عليه وسلم: (أمرت أن أقرأ القرآن على سبعة أحرف) بمعزل... لأنه معلوم أنه لا حرف من حروف القرآن مما اختلفت القراءة في قراءته بهذا المعنى يوجب المراء به كفراً، وقد أوجب صلى الله عليه وسلم بالمراء فيه الكفر».

\section{القول في معنى: (أنزل القرآن من سبعة أبواب الجنة)}
بعد أن فرغ \rawi{الطبري} من بيان معنى الأحرف السبعة، انتقل لبيان معنى الحديث الآخر الذي فيه ذكر الأبواب.

\isnad{قال أبو جعفر:} «اختلف النقلة في ألفاظ الخبر بذلك... فروي عن \rawi{ابن مسعود} عن النبي صلى الله عليه وسلم أنه قال: \begin{hadith}كان الكتاب الأول نزل من باب واحد وعلى حرف واحد، ونزل القرآن من سبعة أبواب وعلى سبعة أحرف: زاجر، وأمر، وحلال، وحرام، ومحكم، ومتشابه، وأمثال\end{hadith}».

وخلاصة رأيه أن \textbf{الأحرف السبعة} تتعلق باختلاف الألفاظ مع اتفاق المعنى (لغات)، بينما \textbf{الأبواب السبعة} تتعلق باختلاف الأوجه والموضوعات التي حواها القرآن (أمر، نهي، حلال، حرام...). فكل وجه من هذه الأوجه هو باب من أبواب الجنة لمن عمل به.

\begin{fawaid}[إشكالات على رأي الطبري]
ذكر الدكتور \rawi{مساعد الطيار} في شرحه أن قول \rawi{الطبري} بأن الصحابة تركوا ستة أحرف من القرآن المنزل هو قول مشكل جداً؛ لأن الأصل أن كل حرف منه قرآن، وما كان قرآناً فلا يجوز لأحد تركه أو إلغاؤه دون نص صريح. كما أن قياسه على التخيير في الكفارة لا يستقيم؛ فالأمة لا يجوز لها الإجماع على ترك كفارتين والإبقاء على واحدة فقط.
\end{fawaid}

بهذا نكون قد انتهينا من هذا الباب المهم من مقدمة الإمام \rawi{الطبري}، وهو من أدق وأصعب مباحث علوم القرآن.